\documentclass[11pt]{article}

\usepackage[T1]{fontenc}
\usepackage[utf8]{inputenc}
\usepackage{lmodern}
\usepackage[margin=1in]{geometry}
\usepackage{microtype}
\usepackage{amsmath,amssymb,mathtools,bm}
\usepackage{booktabs,longtable,tabularx,array}
\usepackage{enumitem}
\usepackage{xcolor}
\usepackage{hyperref}

\hypersetup{
  colorlinks=true,
  linkcolor=blue!55!black,
  urlcolor=blue!55!black,
  pdftitle={Mathematical Guide to the KGParadis Eshelby Implementations}
}

\setcounter{tocdepth}{2}
\setlist{nosep,leftmargin=*}
\allowdisplaybreaks

\newcolumntype{Y}{>{\raggedright\arraybackslash}X}
\newcolumntype{P}[1]{>{\raggedright\arraybackslash}p{#1}}
\newcommand{\vect}[1]{\bm{#1}}
\newcommand{\abs}[1]{\left\lvert #1\right\rvert}
\newcommand{\norm}[1]{\left\lVert #1\right\rVert}
\newcommand{\dd}{\,\mathrm{d}}
\newcommand{\code}[1]{\texttt{\detokenize{#1}}}
\newcommand{\R}{\mathbb{R}}

\title{Mathematical Guide to the KGParadis\\
       \texttt{src/Eshelby*} Implementations}
\author{Source-level technical documentation}
\date{Code snapshot reviewed 7 August 2026}

\begin{document}

\maketitle

\begin{abstract}
This document explains every source file matching \path{src/Eshelby*} in the
reviewed KGParadis tree.  It gives the equations represented by the packed
arrays, relates each loop and matrix operation to a Cartesian multi-index
formula, and derives the direct segment--inclusion nodal force.  Equations
labelled ``literal implementation'' describe what the present source computes;
where that differs from the recentering identity or from another code path, the
difference is stated explicitly rather than silently corrected.
\end{abstract}

\tableofcontents
\newpage

\section{Scope and computational architecture}

There are eight matching files.  Seven implement the Eshelby-specific fast
multipole method (FMM), while one evaluates a nearby inclusion's force on a
straight dislocation segment by analytic line integration.

\begin{center}
\small
\begin{tabularx}{\textwidth}{@{}P{0.35\textwidth}Y P{0.105\textwidth}@{}}
\toprule
File & Role & Guard \\
\midrule
\path{src/EshelbyBuildMatrix.cc} &
  Reconstructs a rectangular contraction block from a packed fully symmetric
  tensor. & none \\
\path{src/EshelbyDerofR.cc} &
  Generates all unique Cartesian derivatives of the inverse-distance kernel
  needed by M2L. & none \\
\path{src/EshelbyEvalTaylor.cc} &
  Evaluates a local Taylor expansion at a point and returns a symmetric stress
  matrix. & \code{ESHELBY} \\
\path{src/EshelbyForce.cc} &
  Integrates the inclusion-induced line-force density against the two linear
  nodal shape functions. & both \\
\path{src/EshelbyMakeEta.cc} &
  Forms raw monomial multipole moments for one spherical inclusion; it is also
  reused to form displacement monomials. & none \\
\path{src/EshelbyMakeTaylor.cc} &
  Performs the multipole-to-local (M2L) operation. & \code{ESHELBY} \\
\path{src/EshelbyTranslateEta.cc} &
  Performs a hard-coded multipole-to-multipole (M2M) binomial translation
  through order six. & \code{ESHELBY} \\
\path{src/EshelbyTranslateTaylor.cc} &
  Performs the local-to-local (L2L) translation used in the downward FMM pass.
  & \code{ESHELBY} \\
\bottomrule
\end{tabularx}
\end{center}

Here ``both'' means that \path{src/EshelbyForce.cc} is nested inside
\code{ESHELBY} and \code{ESHELBYFORCE}.  The far-field sequence is
\[
\boxed{
\text{inclusion}
\xrightarrow{\text{MakeEta}} M
\xrightarrow{\text{TranslateEta}} M'
\xrightarrow[\text{BuildMatrix}]{\text{DerofR + MakeTaylor}} L
\xrightarrow{\text{TranslateTaylor}} L'
\xrightarrow{\text{EvalTaylor}} \vect{\sigma}(\vect{x}) .
}
\]
The direct path \path{src/EshelbyForce.cc} bypasses these expansions and
returns the two endpoint forces directly.

The global setup limits the sum of multipole and Taylor orders to
\[
p+t\leq \code{MAX_ORDER_ESHELBY}=6
\]
in \path{include/FM.h}:14 and \path{src/FMComm.cc}:1464--1485.  This is
essential because the derivative and translation tables stop at that order.

\section{Shared mathematical representation}

\subsection{Multi-indices and packed symmetric tensors}

For a Cartesian multi-index
\[
\alpha=(\alpha_x,\alpha_y,\alpha_z)\in\mathbb{N}_0^3,
\qquad
\abs{\alpha}=\alpha_x+\alpha_y+\alpha_z,
\]
write
\[
\vect{x}^{\alpha}=x^{\alpha_x}y^{\alpha_y}z^{\alpha_z},
\qquad
\alpha!=\alpha_x!\alpha_y!\alpha_z!,
\qquad
\binom{\alpha}{\beta}
=\prod_{q\in\{x,y,z\}}\binom{\alpha_q}{\beta_q}.
\]

The shared table at \path{src/FMComm.cc}:78 is
\[
\code{sequence}[n]=S_n=\binom{n+2}{3}.
\]
It is the number of unique components in all symmetric ranks below (n).
Consequently, a symmetric rank-(n) tensor has
\[
N_n=S_{n+1}-S_n=\binom{n+2}{2}
\]
unique entries and occupies the half-open global interval
\([S_n,S_{n+1})\).

For \(\alpha=(a,b,c)\), \(a+b+c=n\), define \(u=b+c\).  Its local and global
packed indices are
\begin{equation}
\ell_n(a,b,c)=\frac{u(u+1)}{2}+c,
\qquad
I(\alpha)=S_n+\ell_n(a,b,c).
\label{eq:packed-index}
\end{equation}
Thus the order is
\[
\begin{array}{@{}cl@{}}
n&\text{packed components}\ \\
0&1\\
1&x,\ y,\ z\\
2&x^2,\ xy,\ xz,\ y^2,\ yz,\ z^2\\
3&x^3,\ x^2y,\ x^2z,\ xy^2,\ xyz,\ xz^2,
   y^3,\ y^2z,\ yz^2,\ z^3.
\end{array}
\]
The comments made of digits in the source use \(1=x\), \(2=y\), and \(3=z\).
For example, \code{1 1 2 3} means the \(x^2yz\) component.

Moment vectors through degree \(p\) contain \(S_{p+1}\) entries.  Local Taylor
vectors omit ranks zero and one and store ranks two through \(t+2\); their
local packed index and total length are
\begin{equation}
J(\alpha)=I(\alpha)-4,
\qquad
2\leq\abs{\alpha}\leq t+2,
\qquad
N_L=S_{t+3}-4.
\label{eq:taylor-layout}
\end{equation}
In particular, the first six Taylor entries are ordered
\[
(xx,xy,xz,yy,yz,zz).
\]

\subsection{Factorial weights}

The raw monomial and moment arrays do not include factorial denominators.
The array \code{multFactor}, initialized in \path{src/FMComm.cc}:149--255,
contains
\begin{equation}
\code{multFactor}[I(\alpha)-S_2]=\frac{1}{\alpha!},
\qquad 2\leq\abs{\alpha}\leq6.
\label{eq:multfactor}
\end{equation}
Degrees zero and one need no stored factor because their factorials are one.
The M2L, L2L, and evaluation routines apply these factors immediately before a
matrix--vector product.

\subsection{The far-field kernel}

Let \(r=\norm{\vect{r}}\).  The FMM is built from the symmetric tensor
\begin{equation}
K_{ij}(\vect{r})
=\frac{\delta_{ij}}{r^3}-\frac{3r_i r_j}{r^5}
=-\partial_i\partial_j\frac{1}{r}.
\label{eq:kernel}
\end{equation}
For a spherical inclusion of radius (a) and volumetric strain
\[
E_v=\varepsilon_{xx}+\varepsilon_{yy}+\varepsilon_{zz},
\]
the scalar source strength and isotropic material factor used by the FMM are
\begin{equation}
q=E_v a^3,
\qquad
A(\mu,\nu)=\frac{\mu(1+\nu)}{1-2\nu}.
\label{eq:source-strength}
\end{equation}
The monopole field represented by that path is therefore
\begin{equation}
\sigma_{ij}^{\mathrm{FMM}}(\vect{r})
=AqK_{ij}(\vect{r})
=A E_v a^3
 \left(\frac{\delta_{ij}}{r^3}
       -\frac{3r_i r_j}{r^5}\right).
\label{eq:fmm-monopole}
\end{equation}
No factor \(4\pi/3\) is stored: \(E_va^3\), not geometric volume times strain,
is the source normalization in this implementation.

\section{\texttt{EshelbyMakeEta.cc}: particle-to-multipole moments}

The implementation is at \path{src/EshelbyMakeEta.cc}:35--166.  With inclusion
position \(\vect{x}_s\), expansion center \(\vect{c}\), radius \(a=R_0\), and
requested order \(p\), it first forms
\begin{equation}
\vect{d}=\vect{c}-\vect{x}_s,
\qquad q=E_vR_0^3.
\label{eq:makeeta-d}
\end{equation}
The assignments at lines 56--159 enumerate every raw monomial through degree
six.  The final scaling loop at lines 161--163 makes the returned coefficient
for every \(\abs{\alpha}\leq p\)
\begin{equation}
\boxed{M_\alpha=q\vect{d}^{\alpha}.}
\label{eq:makeeta}
\end{equation}
For example,
\[
\begin{aligned}
(M_0,M_x,M_y,M_z)&=q(1,d_x,d_y,d_z),\\
(M_{xx},M_{xy},M_{xz},M_{yy},M_{yz},M_{zz})
&=q(d_x^2,d_xd_y,d_xd_z,d_y^2,d_yd_z,d_z^2).
\end{aligned}
\]
There is deliberately no \(1/\alpha!\) in \eqref{eq:makeeta}; callers apply
\eqref{eq:multfactor} when a Taylor theorem requires it.

For many inclusions in one leaf cell, \path{src/CellCharge.cc}:613--650 sums
these coefficients, so the cell moment is
\begin{equation}
M_\alpha=\sum_s E_{v,s}a_s^3(\vect{c}-\vect{x}_s)^\alpha.
\label{eq:cell-moment}
\end{equation}
That call site also rejects unequal radii for the FMM.  The file's contract is
\(0\leq p\leq6\); the function itself does not validate the order.

\section{\texttt{EshelbyDerofR.cc}: derivatives of the kernel}

The routine at \path{src/EshelbyDerofR.cc}:22--856 accepts a nonzero separation
vector \(\vect{x}\) and fills all distinct components of
\begin{equation}
D_{i_1\ldots i_n}(\vect{x})
=-\partial_{i_1}\cdots\partial_{i_n}\frac{1}{r},
\qquad 2\leq n\leq\code{order}+2.
\label{eq:derofr-definition}
\end{equation}
Rank two is exactly \(K_{ij}\) from \eqref{eq:kernel}.  A compact formula
covering all hard-coded entries is
\begin{equation}
\begin{split}
D_{i_1\ldots i_n}
={}&\sum_{p=0}^{\lfloor n/2\rfloor}
(-1)^{n-p+1}(2n-2p-1)!!\,
r^{-(2n-2p+1)}\\
&\quad\times
\sum_{\mathcal P_p}
\left(\prod_{\{a,b\}\in\mathcal P_p}\delta_{i_ai_b}\right)
\left(\prod_{c\notin V(\mathcal P_p)}x_{i_c}\right),
\end{split}
\label{eq:cartesian-derivative}
\end{equation}
where \(\mathcal P_p\) ranges over all choices of \(p\) disjoint pairs of
index positions, and \(V(\mathcal P_p)\) is the set of paired positions.

The first two nontrivial ranks illustrate the coefficients encoded by the
source:
\begin{align}
D_{ijk}
={}&-\frac{3}{r^5}
 (\delta_{ij}x_k+\delta_{ik}x_j+\delta_{jk}x_i)
 +\frac{15x_ix_jx_k}{r^7},
\label{eq:rank3}\displaybreak[0]\\
D_{ijkl}
={}&-\frac{3}{r^5}
 (\delta_{ij}\delta_{kl}+\delta_{ik}\delta_{jl}
  +\delta_{il}\delta_{jk})
 +\frac{15}{r^7}\operatorname{Sym}_{6}(\delta_{ij}x_kx_l)
 -\frac{105x_ix_jx_kx_l}{r^9}.
\label{eq:rank4}
\end{align}
Here \(\operatorname{Sym}_{6}\) denotes the six distinct placements of one
Kronecker delta and two coordinates.

The inverse-radius cache obeys
\begin{equation}
\code{dinv}[0]=r^{-1},
\qquad
\code{dinv}[m]=(2m-1)!!r^{-(2m+1)},
\qquad 1\leq m\leq8,
\label{eq:dinv}
\end{equation}
and \code{xv} stores the monomials at index \(I(\alpha)\).  Since
\code{drList} omits ranks zero and one,
\begin{equation}
\boxed{
\code{drList}[S_n-4+\ell_n(\alpha)]=D_\alpha,
\qquad n=\abs{\alpha}\geq2.
}
\label{eq:drlist-layout}
\end{equation}

\begin{center}
\begin{tabular}{@{}rrrr@{}}
\toprule
Condition & New rank & \code{drList} indices & Source lines \\
\midrule
always & 2 & 0--5 & 57--62 \\
\code{order > 0} & 3 & 6--15 & 64--89 \\
\code{order > 1} & 4 & 16--30 & 92--162 \\
\code{order > 2} & 5 & 31--51 & 164--259 \\
\code{order > 3} & 6 & 52--79 & 262--399 \\
\code{order > 4} & 7 & 80--115 & 402--592 \\
\code{order > 5} & 8 & 116--160 & 594--851 \\
\bottomrule
\end{tabular}
\end{center}

The required length is
\[
S_{\code{order}+3}-4=\binom{\code{order}+5}{3}-4.
\]
An independent symbolic check of all 161 hard-coded components through rank
eight found agreement with \eqref{eq:derofr-definition} to floating-point
roundoff at a generic nonsymmetric point.  There is no \(r=0\) guard; separated
FMM cells are an input precondition.

\section{\texttt{EshelbyBuildMatrix.cc}: unpacking contraction blocks}

This utility is implemented at \path{src/EshelbyBuildMatrix.cc}:40--87.  For
arguments called Taylor order \(t\) and multipole order \(m\), its active matrix
has
\begin{equation}
n_r=N_{t+2}=\binom{t+4}{2},
\qquad
n_c=N_m=\binom{m+2}{2}.
\label{eq:buildmatrix-size}
\end{equation}
Rows correspond to \(\abs{\alpha}=t+2\), columns to
\(\abs{\beta}=m\), and the mathematical operation is simply
\begin{equation}
\boxed{B_{\alpha\beta}=D_{\alpha+\beta}.}
\label{eq:buildmatrix}
\end{equation}

More explicitly, the row loops at lines 69--83 enumerate
\[
\alpha=(t+2-i,i-j,j),
\quad
r_{\mathrm{row}}=\frac{i(i+1)}2+j,
\quad 0\leq j\leq i\leq t+2,
\]
and the column loops enumerate
\[
\beta=(m-k,k-l,l),
\quad
c_{\mathrm{col}}=\frac{k(k+1)}2+l,
\quad 0\leq l\leq k\leq m.
\]
The source's running offset simplifies to
\[
\code{drOffset}
=\frac{(i+k)(i+k+1)}2+j+l
=\ell_{t+m+2}(\alpha+\beta),
\]
which proves \eqref{eq:buildmatrix}.  Storage is row-major with caller-provided
leading dimension \code{mCols}:
\[
\code{matrix}[r\,\code{mCols}+c].
\]
Only the active upper-left rectangle is cleared; padding in an oversized work
array is left unchanged.  Although the argument is named \code{drList}, the
routine is generic for this packed layout: the Taylor evaluation and L2L paths
also pass a packed Taylor tensor to it.  It applies neither a factorial nor a
material factor.

\section{\texttt{EshelbyMakeTaylor.cc}: multipole-to-local conversion}

The M2L routine is at \path{src/EshelbyMakeTaylor.cc}:30--124.  Let
\(\vect{c}_s\) be the source-cell center, \(\vect{c}_t\) the target-cell center,
and
\[
\vect{R}=\vect{c}_t-\vect{c}_s.
\]
This direction is established by the production call at
\path{src/CellCharge.cc}:284--326.  For every local derivative multi-index
\(\alpha\), \(\abs{\alpha}\leq t\), the generated Taylor tensor is
\begin{equation}
\boxed{
L_{ij,\alpha}
=A(\mu,\nu)
 \sum_{\abs{\beta}\leq p}
 D_{ij,\alpha+\beta}(\vect{R})\frac{M_\beta}{\beta!}.
}
\label{eq:m2l}
\end{equation}
All free indices are symmetric, so \(L_{ij,\alpha}\) is stored as one packed
tensor component of rank \(\abs{\alpha}+2\), using \eqref{eq:taylor-layout}.

The source realizes \eqref{eq:m2l} as follows.
\begin{enumerate}
\item Lines 43--48 allocate
  \(S_{p+t+3}-4\) entries and call \code{EshelbyDerofR} through derivative
  rank \(p+t+2\).
\item Lines 50--86 assemble every \((\abs{\alpha}+2,\abs{\beta})\) block with
  \code{EshelbyBuildMatrix}.  The complete rectangular matrix has
  \[
  (S_{t+3}-4)\ \text{rows and}\ S_{p+1}\ \text{columns}.
  \]
\item Lines 91--105 copy \(M_\beta\) and multiply degrees two and above by
  \(1/\beta!\).  The \code{multFactor} table must already have been
  initialized by \code{EshelbyInitFMArray}.
\item Lines 109--114 perform the matrix--vector product and multiply every
  output by
  \[
  A(\mu,\nu)=\frac{\mu(1+\nu)}{1-2\nu}.
  \]
\end{enumerate}
The separation must be nonzero and \(p+t\leq6\).  The material factor is
singular at \(\nu=1/2\).

\section{\texttt{EshelbyTranslateEta.cc}: multipole translation}

\subsection{The binomial contraction represented by the table}

The M2M implementation is at \path{src/EshelbyTranslateEta.cc}:27--649.  It
first calls
\[
\code{EshelbyMakeEta}(1,1,\code{newCenter},\code{oldCenter},p,\ldots),
\]
so, by \eqref{eq:makeeta-d}, its literal displacement is
\begin{equation}
\vect{\delta}=\vect{c}_{\mathrm{old}}-\vect{c}_{\mathrm{new}}.
\label{eq:literal-shift}
\end{equation}
The outer product at lines 42--54 constructs
\[
\code{mat}_{\beta\gamma}=M_\beta\vect{\delta}^{\gamma}.
\]
The long rank-by-rank table then implements the symmetric-tensor binomial
convolution
\begin{equation}
\boxed{
\widetilde M_\alpha
=\sum_{0\leq\beta\leq\alpha}
  \binom{\alpha}{\beta}
  M_\beta\vect{\delta}^{\alpha-\beta},
\qquad \abs{\alpha}\leq p,
}
\label{eq:m2m-literal}
\end{equation}
apart from the single table entry discussed below.  The implemented blocks are
rank zero and one at lines 56--64, rank two at 67--79, rank three at 82--125,
rank four at 130--217, rank five at 220--373, and rank six at 377--646.

For example, the rank-two and pure rank-three formulas are
\begin{align*}
\widetilde M_{xx}
&=M_{xx}+2M_x\delta_x+M_0\delta_x^2,\\
\widetilde M_{xxx}
&=M_{xxx}+3M_{xx}\delta_x+3M_x\delta_x^2+M_0\delta_x^3.
\end{align*}
The explicit coefficients \(2,3,6,10,15,20,\ldots\) are the multi-index
binomial coefficients in \eqref{eq:m2m-literal}.

\subsection{Literal direction versus the recentering identity}

Moments constructed by \code{EshelbyMakeEta} satisfy
\[
M_\alpha^{\mathrm{old}}
=\sum_s q_s(\vect{c}_{\mathrm{old}}-\vect{x}_s)^\alpha.
\]
Direct construction at the new center therefore requires
\begin{equation}
M_\alpha^{\mathrm{new}}
=\sum_{\beta\leq\alpha}\binom{\alpha}{\beta}
 M_\beta^{\mathrm{old}}
 \vect{h}^{\alpha-\beta},
\qquad
\vect{h}=\vect{c}_{\mathrm{new}}-\vect{c}_{\mathrm{old}}.
\label{eq:m2m-intended}
\end{equation}
The implementation uses \(\vect{\delta}=-\vect{h}\), reversing all terms
with odd shift degree.  This is not just a naming ambiguity: the upward-pass
call at \path{src/FMComm.cc}:2282--2294 passes the parent as the new center and
the child as the old center, consistent with the function comment.  At first
order, the distinction is already visible:
\[
\underbrace{M_x^{\mathrm{old}}+(c_{\mathrm{new},x}-c_{\mathrm{old},x})M_0}
_{\text{direct recentering}}
\quad\hbox{versus}\quad
\underbrace{M_x^{\mathrm{old}}+(c_{\mathrm{old},x}-c_{\mathrm{new},x})M_0}
_{\text{literal source}}.
\]

\subsection{Table audit and resource behavior}

An equation-by-equation audit of all 80 explicit rank-two through rank-six
updates found one index mismatch and no other coefficient/index mismatch.
For \code{ep[41]}, the \(x^2y^3\) component, line 268 contains
\[
\code{mat}(xx,yyy)+\code{mat}(yyy,xz),
\]
where the symmetric second term required by \eqref{eq:m2m-literal} is
\[
\code{mat}(xx,yyy)+\code{mat}(yyy,xx).
\]
In source indices, \code{ELEM(16,6,nCols)} would have to be
\code{ELEM(16,4,nCols)}.  As written, this one term combines to the wrong
multi-index.

The arrays \code{etaTrans} and \code{mat} allocated at lines 37 and 47 are not
freed before the return at line 648, so every call leaks both buffers.  Input
and output may otherwise alias safely because the full outer product is built
before \code{ep} is written.

\section{\texttt{EshelbyTranslateTaylor.cc}: local translation}

The L2L routine is at \path{src/EshelbyTranslateTaylor.cc}:25--126.  For output
derivative order \(i=0,\ldots,t\) and displacement degree
\(j=0,\ldots,t-i\), it asks \code{EshelbyBuildMatrix} for
\begin{equation}
B^{(i,j)}_{\alpha\beta}=L_{\alpha+\beta},
\qquad \abs{\alpha}=i+2,\quad \abs{\beta}=j.
\label{eq:l2l-block}
\end{equation}
The row and column block offsets at lines 72--84 are \(S_{i+2}-4\) and \(S_j\),
respectively.  Lines 105--117 then multiply each disjoint row block by all
admissible displacement monomials.  With correctly normalized monomials, the
literal displacement from \eqref{eq:literal-shift} gives
\begin{equation}
\boxed{
\widetilde L_\alpha
=\sum_{\abs{\beta}\leq t-i}
 L_{\alpha+\beta}\frac{\vect{\delta}^{\beta}}{\beta!},
\qquad \abs{\alpha}=i+2.
}
\label{eq:l2l-literal}
\end{equation}

On the other hand, \code{EshelbyEvalTaylor} uses powers of
\(\vect{x}-\vect{c}\).  The mathematical recentering identity for a polynomial
from \(\vect{c}_{\mathrm{old}}\) to \(\vect{c}_{\mathrm{new}}\) is therefore
\begin{equation}
L_\alpha^{\mathrm{new}}
=\sum_{\abs{\beta}\leq t-i}
 L_{\alpha+\beta}^{\mathrm{old}}
 \frac{\vect{h}^{\beta}}{\beta!},
\qquad \vect{h}=\vect{c}_{\mathrm{new}}-\vect{c}_{\mathrm{old}}.
\label{eq:l2l-intended}
\end{equation}
The file again uses \(\vect{\delta}=-\vect{h}\).  The downward call at
\path{src/FMComm.cc}:2074--2125 passes the parent center as old and child center
as new, so odd shift-degree terms have the opposite sign from
\eqref{eq:l2l-intended} under the documented evaluation convention.

There is also a precise endpoint error in factorial scaling.  Lines 51--55 set
the last valid index to \(S_{t+1}-1\) but loop while the index is strictly less
than that value.  For every \(t\geq2\), all needed monomials are divided by
their factorial except the final packed component:
\begin{equation}
q_\beta=\frac{\vect{\delta}^{\beta}}{\beta!}
\quad\text{except}\quad
q_{(0,0,t)}=\delta_z^t.
\label{eq:l2l-factor-bug}
\end{equation}
Thus only the output rank-two block can receive the highest-degree error, with
excess
\[
\left(1-\frac{1}{t!}\right)
L_{\alpha+(0,0,t)}\delta_z^t,
\qquad \abs{\alpha}=2.
\]
The temporary matrix is freed at lines 89--90, and the two remaining buffers
are freed at lines 119--123.

\section{\texttt{EshelbyEvalTaylor.cc}: local expansion at a point}

The implementation is at \path{src/EshelbyEvalTaylor.cc}:24--100.  Let
\[
\vect{y}=\vect{x}-\vect{c}_t.
\]
The call at line 41 deliberately reverses the apparently descriptive
arguments to \code{EshelbyMakeEta}:
\[
\code{position}=\vect{c}_t,
\qquad
\code{center}=\vect{x},
\]
which correctly produces raw powers \(\vect{y}^\alpha\).  Lines 69--75 apply
every \(1/\alpha!\) factor, including the final component.  The evaluated
stress is
\begin{equation}
\boxed{
\sigma_{ij}(\vect{x})
=\sum_{\abs{\alpha}\leq t}
 L_{ij,\alpha}\frac{\vect{y}^{\alpha}}{\alpha!}.
}
\label{eq:taylor-eval}
\end{equation}

For each degree, \code{EshelbyBuildMatrix(taylorSeries,0,i,...)} reconstructs
the six stress rows \(ij\in(xx,xy,xz,yy,yz,zz)\) from the packed fully
symmetric local tensor.  After the matrix--vector product, lines 82--90 return
\begin{equation}
\vect{\sigma}=
\begin{pmatrix}
\sigma_{xx}&\sigma_{xy}&\sigma_{xz}\\
\sigma_{xy}&\sigma_{yy}&\sigma_{yz}\\
\sigma_{xz}&\sigma_{yz}&\sigma_{zz}
\end{pmatrix}.
\label{eq:stress-unpack}
\end{equation}
The allocated matrix and output vector have \(S_{t+3}-4\) rows even though
only their first six rows are populated or consumed; this is extra work and
storage, not a change to \eqref{eq:taylor-eval}.

\section{\texttt{EshelbyForce.cc}: analytic endpoint force}

\subsection{Inputs and model restrictions}

The helpers are at \path{src/EshelbyForce.cc}:31--103 and the public routine at
lines 134--287.  It uses the inclusion center \(\vect{c}\), three radius
entries, the six-component symmetric strain vector, two segment endpoints,
and the Burgers vector \(\vect{b}\).  The strain layout is
\[
(\varepsilon_{xx},\varepsilon_{yy},\varepsilon_{zz},
  \varepsilon_{yz},\varepsilon_{xz},\varepsilon_{xy}),
\]
but the routine uses only
\[
E_v=\varepsilon_{xx}+\varepsilon_{yy}+\varepsilon_{zz}.
\]
The source comment assumes the three diagonal entries are equal, although the
code checks neither equality nor absence of deviatoric diagonal strain.  All
shear entries are ignored.

Although the inclusion type stores three semi-principal axes, this routine sets
\begin{equation}
a=\max(a_x,a_y,a_z)
\label{eq:max-radius}
\end{equation}
and performs a spherical intersection.  It does not receive or use the
ellipsoid orientation.  Matrix properties are \((\mu_1,\nu_1)\), while
inclusion properties are \((\mu_2,\nu_2)\); initialization makes material 2
equal to material 1 if no second values were supplied.

\subsection{Segment geometry and sphere intersection}

Let the endpoints be \(\vect{x}_1,\vect{x}_2\), and define
\begin{equation}
\vect{\ell}=\vect{x}_2-\vect{x}_1,
\qquad L=\norm{\vect{\ell}},
\qquad \vect{t}=\frac{\vect{\ell}}{L}.
\label{eq:segment-tangent}
\end{equation}
With \(\vect{R}_i=\vect{x}_i-\vect{c}\), the axial coordinates and closest
approach vector are
\begin{equation}
s_0=\vect{R}_1\mathbin{\cdot}\vect{t},
\qquad
s_1=\vect{R}_2\mathbin{\cdot}\vect{t}=s_0+L,
\qquad
\vect{n}=\vect{R}_1-s_0\vect{t},
\qquad d^2=\norm{\vect{n}}^2.
\label{eq:line-coordinates}
\end{equation}
Hence a segment point and its distance from the inclusion center are
\begin{equation}
\vect{r}(s)=\vect{n}+s\vect{t},
\qquad
\rho(s)=\sqrt{d^2+s^2},
\qquad s\in[s_0,s_1].
\label{eq:line-parameterization}
\end{equation}
The infinite line intersects the effective sphere when
\(\Delta=a^2-d^2>0\), at
\begin{equation}
s_-=-\sqrt{\Delta},
\qquad s_+=+\sqrt{\Delta}.
\label{eq:sphere-intersection}
\end{equation}

Lines 209--244 divide the segment into exterior set \(\mathcal O\) and
interior set \(\mathcal I\):
\begin{center}
\begin{tabularx}{\textwidth}{@{}Y Y Y@{}}
\toprule
Case & \(\mathcal O\) & \(\mathcal I\) \\
\midrule
line miss or no segment overlap & \([s_0,s_1]\) & \(\varnothing\) \\
starts inside, exits & \([s_+,s_1]\) & \([s_0,s_+]\) \\
enters, ends inside & \([s_0,s_-]\) & \([s_-,s_1]\) \\
wholly inside & \(\varnothing\) & \([s_0,s_1]\) \\
passes through & \([s_0,s_-]\cup[s_+,s_1]\) & \([s_-,s_+]\) \\
\bottomrule
\end{tabularx}
\end{center}
For two exterior pieces, their five integral vectors are added componentwise.

\subsection{Interior and exterior primitive integrals}

For an interior interval \([p,q]\), \code{EshelbySpecialIntegrals} returns
\begin{equation}
S_0=\int_p^q1\dd s=q-p,
\qquad
S_1=\int_p^qs\dd s=\frac{q^2-p^2}{2}.
\label{eq:inside-integrals}
\end{equation}

For an exterior interval, \code{EshelbyIntegrals} returns endpoint differences
of the following antiderivatives:
\begin{align}
\Phi_{03}(s)&=\frac{s}{d^2\rho},
&\Phi_{13}(s)&=-\frac{1}{\rho},\nonumber\\
\Phi_{05}(s)&=\frac{2\Phi_{03}(s)+s/\rho^3}{d^2},
&\Phi_{15}(s)&=-\frac{1}{\rho^3},\nonumber\\
\Phi_{25}(s)&=\Phi_{03}(s)-\frac{s}{\rho^3}.
\label{eq:force-primitives}
\end{align}
Thus the array \code{wInt} has the exact meanings
\begin{equation}
\begin{aligned}
W_0&=\int_{\mathcal O}\rho^{-3}\dd s,
&W_1&=\int_{\mathcal O}s\rho^{-3}\dd s,\\
W_2&=3\int_{\mathcal O}\rho^{-5}\dd s,
&W_3&=3\int_{\mathcal O}s\rho^{-5}\dd s,\\
W_4&=3\int_{\mathcal O}s^2\rho^{-5}\dd s.
\end{aligned}
\label{eq:wint}
\end{equation}
The factors of three in the last three entries are part of the array
definition.

If
\[
d^2<10^{-8}(p^2+q^2),
\]
the source uses the collinear limit
\[
\Phi_{03}=-\frac{s}{2\rho^3},
\qquad \Phi_{13}=-\rho^{-1},
\qquad \Phi_{05}=\Phi_{15}=\Phi_{25}=0.
\]
At exactly \(d=0\), the zeroed terms multiply
\(\vect{n}\times\vect{t}=\vect{0}\); for small nonzero \(d\), this is an
approximation.  The code computes \(1/d^2\) before taking this branch, so an
exactly collinear case still produces an unused floating-point infinity.

\subsection{Shape functions and endpoint moments}

The linear segment shape functions are
\begin{equation}
N_1(s)=\frac{s_1-s}{L},
\qquad
N_2(s)=\frac{s-s_0}{L}.
\label{eq:shape-functions}
\end{equation}
For endpoint 2, lines 258--261 form
\begin{align}
F_{03}^{(2)}&=W_1-s_0W_0,
&F_{05}^{(2)}&=W_3-s_0W_2,\nonumber\\
F_{15}^{(2)}&=W_4-s_0W_3,
&F_{\mathrm{SP}}^{(2)}&=S_1-s_0S_0.
\label{eq:endpoint2-moments}
\end{align}
For endpoint 1, lines 275--278 form
\begin{align}
F_{03}^{(1)}&=s_1W_0-W_1,
&F_{05}^{(1)}&=s_1W_2-W_3,\nonumber\\
F_{15}^{(1)}&=s_1W_3-W_4,
&F_{\mathrm{SP}}^{(1)}&=s_1S_0-S_1.
\label{eq:endpoint1-moments}
\end{align}
For example,
\[
F_{05}^{(2)}=3\int_{\mathcal O}(s-s_0)\rho^{-5}\dd s,
\qquad
F_{\mathrm{SP}}^{(1)}=\int_{\mathcal I}(s_1-s)\dd s.
\]
Division by \(L\) is applied later through the material factors.

\subsection{Vector and material contraction}

The three fixed vectors at lines 246--256 are
\begin{equation}
\vect{A}=\vect{b}\times\vect{t},
\qquad
\vect{B}=(\vect{b}\mathbin{\cdot}\vect{n})(\vect{n}\times\vect{t}),
\qquad
\vect{D}=(\vect{b}\mathbin{\cdot}\vect{t})(\vect{n}\times\vect{t}).
\label{eq:force-vectors}
\end{equation}
Define
\begin{equation}
C_j=\frac{\mu_j(1+\nu_j)}{1-2\nu_j}E_v,
\qquad j\in\{1,2\}.
\label{eq:force-material}
\end{equation}
The exact endpoint equations implemented at lines 270--283 are
\begin{equation}
\boxed{
\vect{f}_i
=-\frac{C_1a^3}{L}
 \left(\vect{A}F_{03}^{(i)}
      +\vect{B}F_{05}^{(i)}
      +\vect{D}F_{15}^{(i)}\right)
 +\frac{2C_2}{L}\vect{A}F_{\mathrm{SP}}^{(i)},
\qquad i=1,2.
}
\label{eq:endpoint-force}
\end{equation}
Equivalently, this is the shape-function integral of the literal line-force
density
\begin{align}
\vect{g}_{\mathrm{out}}(s)
&=-C_1a^3\left[
 \frac{\vect{b}\times\vect{t}}{\rho^3}
 +3\frac{(\vect{b}\mathbin{\cdot}\vect{r})
          (\vect{r}\times\vect{t})}{\rho^5}
 \right],
\label{eq:outside-density}\\
\vect{g}_{\mathrm{in}}(s)
&=2C_2(\vect{b}\times\vect{t}),
\label{eq:inside-density}\\
\vect{f}_i&=\int_{\mathcal O}N_i\vect{g}_{\mathrm{out}}\dd s
            +\int_{\mathcal I}N_i\vect{g}_{\mathrm{in}}\dd s.
\label{eq:nodal-integral}
\end{align}
Both nodal forces are perpendicular to \(\vect{t}\).

Using the standard Peach--Koehler relation
\(\vect{g}=(\vect{\sigma}\vect{b})\times\vect{t}\), the stress fields
algebraically implied by this file are
\begin{equation}
\boxed{
\vect{\sigma}_{\mathrm{out}}^{\mathrm{force\ file}}
=-C_1a^3\left(\frac{\vect{I}}{\rho^3}
              +3\frac{\vect{r}\vect{r}^{T}}{\rho^5}\right),
\qquad
\vect{\sigma}_{\mathrm{in}}^{\mathrm{force\ file}}=2C_2\vect{I}.
}
\label{eq:force-implied-stress}
\end{equation}

\section{Cross-file consistency and implementation cautions}

This section records literal source behavior that is easy to miss when reading
only the comments.

\subsection{Direct-force and FMM fields are not the same kernel}

The FMM monopole \eqref{eq:fmm-monopole} is
\[
\vect{\sigma}_{\mathrm{out}}^{\mathrm{FMM}}
=+C_1a^3\left(\frac{\vect{I}}{\rho^3}
              -3\frac{\vect{r}\vect{r}^T}{\rho^5}\right).
\]
The direct routine implies \eqref{eq:force-implied-stress}, which differs in
the overall isotropic sign and in the relative sign of the radial term.  Its
interior value also has the opposite sign from the disabled reference stress
at \path{src/RemoteSegForces.cc}:290--335, which states
\(-2C_1\vect{I}\).  The material-2 constants can intentionally alter the
interior magnitude, but do not explain these sign differences.  Therefore
near-field direct and far-field FMM inclusion forces are not algebraically
consistent under the Peach--Koehler convention used by
\path{src/RemoteSegForces.cc}:903--905.  A numerical regression test spanning
the near/far cutoff is warranted before choosing which formula to change.

\subsection{Translation discrepancies}

Both translation functions obtain
\(\vect{c}_{\mathrm{old}}-\vect{c}_{\mathrm{new}}\) from
\code{EshelbyMakeEta}, while the moment definition and Taylor evaluation
require \(\vect{c}_{\mathrm{new}}-\vect{c}_{\mathrm{old}}\).  This reverses
all odd shift-degree contributions in M2M and L2L.  In addition:
\begin{itemize}
\item \path{src/EshelbyTranslateEta.cc}:268 has one wrong packed matrix index
  in the \(x^2y^3\) coefficient;
\item that function never frees its two temporary heap arrays;
\item \path{src/EshelbyTranslateTaylor.cc}:51--55 fails to divide the final
  \(z^t\) displacement monomial by \(t!\).
\end{itemize}

\subsection{Input and boundary assumptions}

\begin{itemize}
\item The FMM tables require nonnegative orders with \(p+t\leq6\), but the
  individual functions generally rely on the global setup rather than checking
  these bounds locally.
\item \code{EshelbyDerofR} divides by \(r\), and \code{EshelbyForce} divides by
  segment length \(L\); neither accepts a zero value.
\item All factors containing \(1-2\nu\) assume \(\nu\neq1/2\).
\item The direct-force routine replaces any ellipsoid by a sphere of maximum
  semi-axis, whereas the FMM rejects unequal radii.
\item Exact, strict intersection comparisons at
  \path{src/EshelbyForce.cc}:209--235 make an endpoint exactly on the sphere a
  delicate boundary case.  Certain ``boundary to interior'' configurations
  fall into the final through-going branch and can create a reversed exterior
  interval.
\item The no-overlap condition uses bitwise \code{|} rather than logical
  \code{||}.  Its truth value is the same for ordinary comparison results, but
  it evaluates all three operands.
\end{itemize}

\section{Compact end-to-end equation summary}

For a well-separated source cell and target cell, the nominal mathematical
pipeline represented by the files can be summarized as
\begin{align}
M_\beta^{(s)}
&=\sum_a E_{v,a}a_a^3
  (\vect{c}_s-\vect{x}_a)^\beta,
&&\text{P2M: \code{EshelbyMakeEta}},
\label{eq:summary-p2m}\\
M_\alpha^{(p)}
&=\sum_{\beta\leq\alpha}\binom{\alpha}{\beta}
  M_\beta^{(s)}
  (\vect{c}_p-\vect{c}_s)^{\alpha-\beta},
&&\text{M2M: intended recentering},
\label{eq:summary-m2m}\\
L_{ij,\alpha}
&=A\sum_{\abs{\beta}\leq p}
  D_{ij,\alpha+\beta}(\vect{c}_t-\vect{c}_s)
  \frac{M_\beta}{\beta!},
&&\text{M2L: \code{DerofR/BuildMatrix/MakeTaylor}},
\label{eq:summary-m2l}\\
L_{ij,\alpha}^{(c)}
&=\sum_{\abs{\beta}\leq t-\abs{\alpha}}
  L_{ij,\alpha+\beta}^{(p)}
  \frac{(\vect{c}_c-\vect{c}_p)^\beta}{\beta!},
&&\text{L2L: intended recentering},
\label{eq:summary-l2l}\\
\sigma_{ij}(\vect{x})
&=\sum_{\abs{\alpha}\leq t}
  L_{ij,\alpha}
  \frac{(\vect{x}-\vect{c}_t)^\alpha}{\alpha!},
&&\text{L2P: \code{EshelbyEvalTaylor}}.
\label{eq:summary-l2p}
\end{align}
Equations \eqref{eq:summary-m2m} and \eqref{eq:summary-l2l} are the nominal
recenterings.  The present translation sources use the negatives of the shown
center displacements, as detailed in their respective sections.

\appendix
\section{Source-to-equation index}

\small
\begin{longtable}{@{}P{0.38\textwidth}P{0.22\textwidth}P{0.32\textwidth}@{}}
\toprule
File and lines & Main equations & Principal arrays/outputs \\
\midrule
\endhead
\path{src/EshelbyBuildMatrix.cc}:40--87 &
\eqref{eq:buildmatrix-size}, \eqref{eq:buildmatrix} &
Rectangular row-major block \(B_{\alpha\beta}\). \\
\path{src/EshelbyDerofR.cc}:22--856 &
\eqref{eq:derofr-definition}--\eqref{eq:drlist-layout} &
Packed ranks 2--8 in \code{drList}; monomial cache \code{xv}. \\
\path{src/EshelbyEvalTaylor.cc}:24--100 &
\eqref{eq:taylor-eval}, \eqref{eq:stress-unpack} &
Symmetric \(3\times3\) \code{stressMatrix}. \\
\path{src/EshelbyForce.cc}:31--103, 134--287 &
\eqref{eq:line-parameterization}, \eqref{eq:wint},
\eqref{eq:endpoint-force} &
Five exterior integrals, two interior integrals, endpoint forces
\code{f1}/\code{f2}. \\
\path{src/EshelbyMakeEta.cc}:35--166 &
\eqref{eq:makeeta-d}--\eqref{eq:cell-moment} &
Raw moment/monomial vector \code{etav}. \\
\path{src/EshelbyMakeTaylor.cc}:30--124 &
\eqref{eq:m2l} &
Packed local tensor \code{taylorSeries}. \\
\path{src/EshelbyTranslateEta.cc}:27--649 &
\eqref{eq:m2m-literal}, \eqref{eq:m2m-intended} &
Translated moments \code{ep}; explicit ranks 0--6. \\
\path{src/EshelbyTranslateTaylor.cc}:25--126 &
\eqref{eq:l2l-block}--\eqref{eq:l2l-factor-bug} &
Translated local tensor \code{tPrime}. \\
\bottomrule
\end{longtable}
\normalsize

\section{Notation quick reference}

\begin{tabularx}{\textwidth}{@{}p{0.19\textwidth}Y@{}}
\toprule
Symbol & Meaning \\
\midrule
\(S_n=\binom{n+2}{3}\) & Start index of packed symmetric rank \(n\). \\
\(I(\alpha)\) & Global packed index including ranks zero and one. \\
\(J(\alpha)=I(\alpha)-4\) & Packed Taylor/derivative index, with ranks zero and one omitted. \\
\(M_\alpha\) & Raw, unnormalized multipole moment. \\
\(L_{ij,\alpha}\) & Local Taylor derivative coefficient. \\
\(D_{i_1\ldots i_n}\) & \(-\partial_{i_1}\cdots\partial_{i_n}(1/r)\). \\
\(A(\mu,\nu)\) & \(\mu(1+\nu)/(1-2\nu)\). \\
\(E_v\) & Trace of the stored inclusion strain. \\
\(a\) & Sphere radius; in the direct routine, maximum of the three radii. \\
\(p,t\) & Multipole and Taylor truncation orders. \\
\bottomrule
\end{tabularx}

\end{document}
